\section{Conclusion}
\label{conclusion}

Ce mémoire partait d'une asymétrie~: le Front Office connaît le gain de chaque deal,
quand le Back Office ne connaît que son coût global, jamais celui d'un contrat pris
isolément. Faute d'instrument, le coût opérationnel d'un deal de crédit syndiqué restait
une inconnue --- et avec lui, la rentabilité réelle du contrat. D'où la problématique~:
\emph{comment formaliser la charge de travail des équipes Back Office en une mesure
quantitative fiable au service de la décision~?} La réponse tient en une chaîne de
mesure, construite le long des trois sous-questions posées en introduction.

\paragraph{Mesurer.} L'unité retenue est l'événement opérationnel~: assez fin pour être
chronométré de façon homogène, assez structurant pour que tout deal se recompose comme la
somme de ses événements. La charge unitaire se scinde en deux objets de nature
différente --- un \textit{Event Processing} au temps de booking remarquablement stable
($\bar{t}\approx 20$~minutes, sur environ 500 \textit{clockings}) et un \textit{Aftercare}
variable, qui porte l'essentiel de la complexité. Cette mesure, inspirée de l'étude des
temps mais expurgée du facteur d'allure, a été conduite dans une démarche explicitement
diplomatique et calée hors des pics d'activité, pour peser le moins possible sur les
équipes observées.

\paragraph{Modéliser.} La charge unitaire, valorisée au coût horaire chargé, reconstruit
rétrospectivement le coût d'un deal selon l'équation du TDABC,
$C=(T_{LIQ}+T_{AC})\times N \times w_s$. Pour la prédiction, une architecture à deux étages
sépare ce qui obéit à des logiques distinctes~: un modèle à base d'arbres (Random Forest,
XGBoost) pour le volume d'événements, dont les facteurs interagissent de façon non
additive~; une régression paramétrique pour l'Aftercare, dont la relation aux
caractéristiques du deal reste proportionnelle et interprétable. Les valeurs SHAP
restituent, deal par deal, la lisibilité qu'exige un arbitrage porté devant le Front
Office.

\paragraph{Décider.} La mesure débouche sur trois leviers de pilotage~: cibler les
initiatives de réduction des coûts sur la frange critique des deals, évaluer le
\textit{cost avoidance} réel du nearshoring en comparant le coût d'un même portefeuille
entre sites, et projeter la charge à venir sur le pipeline de deals signés.

Le résultat empirique le plus marquant est la dispersion extrême des coûts~: d'un deal
standard à quelques milliers d'euros à un \textit{jumbo deal} dépassant plusieurs
centaines de milliers, l'écart atteint un rapport de l'ordre de 1 à 60, et 1~\% des deals
concentrent près de 30~\% du volume total d'événements. Cette concentration justifie à
elle seule l'effort de modélisation~: c'est sur une poignée de contrats que se joue
l'essentiel de la charge Back Office.

La portée de ces résultats appelle des réserves assumées. Les temps mesurés sont des
ordres de grandeur robustes, non des chronométrages absolus~: l'effet Hawthorne a été
contenu, non éliminé. Certains déterminants de la charge --- complexité juridique,
qualité de la documentation amont, clauses atypiques --- échappent aux données
disponibles et nourrissent le bruit résiduel. Enfin, un modèle appris sur l'historique
suppose une gouvernance de ré-entraînement pour résister à la dérive des distributions.

Ces limites tracent les prolongements. L'intégration des données de mails (\textit{Hobart})
affinerait l'estimation de l'Aftercare~; des modèles hybrides capteraient conjointement
proportionnalité et interactions~; l'extension à un horizon multi-mois servirait la
planification capacitaire. Surtout, l'irruption de l'IA générative --- qui déplace le
travail humain vers les cas complexes, précisément les plus coûteux --- fait de la mesure
de la charge un préalable à toute automatisation raisonnée~: on n'optimise bien que ce
qu'on a d'abord su mesurer.

En rendant le coût d'un deal comparable au gain qu'en connaît le Front Office, ce travail
dote le Back Office de l'instrument qui lui manquait pour parler le langage de la
rentabilité, deal par deal.
